%% ============================================================================
%%  Closed-Loop LLM Task Planning for a Simulated Mobile Manipulator
%%  IEEE conference (IEEEtran) draft. Compile in Overleaf with pdfLaTeX.
%%
%%  All experimental numbers are the FROZEN-SYSTEM results (runs on/after
%%  2026-06-25). See PAPER_DOSSIER.md for provenance. Do NOT change any number
%%  to a rounder/higher value without checking the logs.
%% ============================================================================
\documentclass[conference]{IEEEtran}
\IEEEoverridecommandlockouts
\usepackage{cite}
\usepackage{amsmath,amssymb,amsfonts}
\usepackage{algorithmic}
\usepackage{graphicx}
\usepackage{textcomp}
\usepackage{xcolor}
\usepackage{booktabs}
\usepackage{url}
\def\BibTeX{{\rm B\kern-.05em{\sc i\kern-.025em b}\kern-.08em T\kern-.1667em\lower.7ex\hbox{E}\kern-.125emX}}

\begin{document}

\title{Closed-Loop Large Language Model Task Planning for a Simulated Mobile Manipulator:\\
An Implementation Study and Honest Reliability Analysis}

\author{\IEEEauthorblockN{Firaol Tesfaye}
\IEEEauthorblockA{\textit{[Department / Institution]} \\
\textit{[University]}\\
firaalx@gmail.com}}

\maketitle

\begin{abstract}
We present a complete, reproducible implementation of a closed-loop large
language model (LLM) task planner for a mobile manipulator, built entirely on a
standard open-source robotics stack: ROS~2 Jazzy, Gazebo Harmonic, Nav2, MoveIt~2,
and YOLOv8. An LLM (Anthropic Claude Haiku~4.5) receives a natural-language
command and a structured scene state, and acts only by invoking a fixed set of
robot \emph{tool services} (navigate, detect, pick, place, drive, \dots); a fresh
scene state is returned after every action, closing the perception--action loop
and enabling the planner to \emph{replan} after failures. We compare this against
a hardcoded baseline controller that executes the same task through the same tool
services but with a fixed, perception-blind action sequence and no recovery. On
the frozen system, single-object pick-and-deliver succeeds in $13/16$ ($81\%$)
of trials, with nearly every success involving autonomous recovery from an initial
out-of-reach grasp; the deterministic baseline succeeds reliably on its tuned
configuration but cannot adapt. Reliability degrades on longer, multi-object
sequences ($2/4$), and an ambiguous ``tidy up the warehouse'' command is
correctly \emph{interpreted} by the LLM but exposes a live-perception limitation
during execution. We report these outcomes honestly, including residual failure
modes, and distill a set of concrete simulation--robotics integration lessons.
\end{abstract}

\begin{IEEEkeywords}
large language models, task planning, mobile manipulation, ROS~2, closed-loop
control, robot learning, human--robot instruction
\end{IEEEkeywords}

% ============================================================================
\section{Introduction}
% ============================================================================
Instructing a robot in natural language---``take the red cube to
dispatch''---and having it decompose, execute, and \emph{recover} without
task-specific scripting is a long-standing goal of service and warehouse
robotics. Large language models (LLMs) have recently made the decomposition step
tractable, but a language model alone cannot move a robot: it must be grounded in
a real perception--action stack, and it must react when the physical world does
not match its plan.

This paper is an \emph{implementation study}. Rather than proposing a new
algorithm, we ask a practical question: \emph{how far does a general-purpose LLM,
wired into a conventional ROS~2 stack as a closed-loop tool-caller, actually get
on a mobile-manipulation task---and where, honestly, does it fail?} We build the
full system (Section~\ref{sec:system}), define a hardcoded baseline that isolates
\emph{decision-making} from \emph{capability} (Section~\ref{sec:baseline}), and
evaluate both on single-object, multi-object, and ambiguous-command scenarios
(Section~\ref{sec:results}).

Our contributions are:
\begin{enumerate}
\item A complete, reproducible closed-loop LLM mobile-manipulation system on
open-source components (ROS~2 Jazzy, Gazebo Harmonic, Nav2 MPPI, MoveIt~2
RRTConnect, YOLOv8), with the LLM restricted to a fixed tool interface.
\item A controlled comparison against a hardcoded baseline that shares the same
tool services, isolating the effect of closed-loop LLM decision-making.
\item An \emph{honest} reliability analysis: $81\%$ single-object success with
autonomous recovery, measured degradation on longer tasks, a characterized
ambiguous-command limitation, and a catalogue of integration bugs and fixes that
others building similar systems will hit.
\end{enumerate}

% ============================================================================
\section{Related Work}
% ============================================================================
LLM-driven robot planning was popularized by approaches that ground language in a
library of robot skills and score feasible next actions, e.g.\ SayCan-style
affordance grounding, and by systems that let an LLM emit code or structured tool
calls against a robot API. Our design follows the tool-calling paradigm: the LLM
never emits motor commands directly, only calls to typed services with validated
arguments. We differ from many demonstrations in two ways. First, we close the
loop with a \emph{structured scene state returned after every action}, so the
planner reacts to real execution results rather than open-loop emitting a plan.
Second, we deliberately foreground \emph{failure analysis}: like the appendices of
several LLM-robotics papers, we treat characterized failure modes as a result, not
an embarrassment. A full related-work treatment is left for the camera-ready.

% ============================================================================
\section{System Architecture}\label{sec:system}
% ============================================================================
\subsection{Platform}
The robot is a differential-drive base carrying a 6-DOF arm ($6$ revolute joints,
each limited to $\pm3.14$~rad), a two-finger gripper, and an RGB-D camera. It is
simulated in Gazebo Sim~8.11 (Harmonic) with the DART physics engine on Ubuntu
22.04/24.04 under ROS~2 Jazzy. Navigation uses Nav2 with the MPPI controller on a
frozen static map (identity $map\rightarrow odom$ transform; no live SLAM during
experiments, for reasons discussed in Section~\ref{sec:lessons}). Arm motion uses
MoveIt~2 (\texttt{MoveGroupInterface}) with the OMPL RRTConnect planner. Object
detection uses Ultralytics YOLOv8n.

The software is organized into seven ROS~2 packages. Per project convention, all
nodes are C++ (rclcpp) except exactly two Python (rclpy) nodes---the YOLO
perception node and the LLM client---reflecting the practical reality that the ML
tooling lives in Python while the real-time control and service layer benefit from
C++.

\subsection{Tool server}
A C++ \emph{tool server} exposes the robot's capabilities as ROS~2 services, each
a self-contained, blocking, verifiable action. The eleven tools available to the
planner are: \texttt{navigate\_to}, \texttt{drive\_distance},
\texttt{detect\_objects}, \texttt{move\_arm\_to\_named} (\texttt{home}/\texttt{ready}),
\texttt{move\_arm\_to\_pose}, \texttt{pick}, \texttt{place},
\texttt{open\_gripper}, \texttt{close\_gripper}, \texttt{get\_state}, and
\texttt{task\_complete}. \texttt{navigate\_to} drives Nav2 to one of four fixed
named locations (\texttt{table}, \texttt{dispatch\_zone}, \texttt{shelf\_area},
\texttt{home}); \texttt{drive\_distance} performs a short blind dead-reckoning
move (capped at $\pm2$~m) for local adjustments; \texttt{pick} performs a
perception-driven grasp using a DetachableJoint attach in simulation.

\subsection{Closed-loop LLM planner}
The planner (Fig.~\ref{fig:loop}) exposes a ROS~2 action \texttt{/execute\_task}.
Given a command, it fetches an initial scene state, starts an LLM conversation
governed by a fixed system prompt, and then loops up to a maximum number of turns
(default $30$). At each turn the LLM either requests a tool call---dispatched
generically by constructing the service request from a schema and polling the
result with a $200$~s deadline---or calls \texttt{task\_complete}. Crucially,
after every \texttt{navigate\_to}, \texttt{drive\_distance}, \texttt{pick}, or
\texttt{place}, a \emph{fresh structured scene state is appended} to the
conversation, giving the LLM up-to-date robot pose and detected-object
coordinates. Repeated identical failures trigger an injected hint to break loops.
Every turn, tool call, and the full conversation are logged to JSON for analysis.

\begin{figure}[t]
\centering
% Replace with the architecture/loop diagram (see PAPER_DOSSIER.md for content).
\framebox[0.95\columnwidth]{\parbox[c][3.2cm][c]{0.9\columnwidth}{\centering
\textit{Fig.~1 placeholder:} command $\rightarrow$ LLM $\rightarrow$ tool call
$\rightarrow$ tool server (Nav2 / MoveIt / YOLO) $\rightarrow$ fresh scene state
$\rightarrow$ back to LLM, until \texttt{task\_complete}.}}
\caption{Closed-loop LLM task-planning architecture. A structured scene state is
returned after every action, enabling replanning.}
\label{fig:loop}
\end{figure}

\subsection{The single most important reliability margin}
The base navigates to a standoff in front of the table, which leaves the cubes at
the very edge of arm reach; the first \texttt{pick} therefore typically fails with
a pre-grasp IK error. A short forward \texttt{drive\_distance(0.30\,m)} ``approach
nudge'' brings the cube into reach. Whether this nudge is issued---and it is not
hardcoded into the LLM path---is the dominant factor in single-object success, and
the LLM learns to issue it from the failure message (Section~\ref{sec:results}).

% ============================================================================
\section{Baseline Controller}\label{sec:baseline}
% ============================================================================
To isolate the contribution of closed-loop LLM decision-making from raw
capability, the baseline executes the \emph{same task through the same tool
services}, differing only in how the next action is chosen. It runs a fixed
sequence---\texttt{navigate\_to(table)} $\rightarrow$
\texttt{drive\_distance(0.30)} $\rightarrow$ pick-at-hardcoded-pose $\rightarrow$
\texttt{navigate\_to(dispatch\_zone)} $\rightarrow$ \texttt{place}---with three
deliberate properties of a traditional automation program: (i)~it picks at a
compile-time-constant coordinate and \emph{never} calls perception; (ii)~it does
not branch on scene state; and (iii)~it has \emph{no recovery}---any failed step
aborts the run. A two-cube mode repeats the sequence for a second object.

% ============================================================================
\section{Experimental Setup}\label{sec:experiments}
% ============================================================================
An automated experiment runner resets the world before each trial (release and
open the gripper, move the arm home, restore cube spawn poses, return the base
home), then issues the command and records success, wall-clock duration, path
length, and---for the LLM---the number of LLM calls and replans. We evaluate four
scenarios (Table~\ref{tab:scenarios}).

\begin{table}[t]
\caption{Evaluation scenarios.}
\label{tab:scenarios}
\centering
\small
\begin{tabular}{@{}llp{4.4cm}@{}}
\toprule
ID & System & Command / purpose \\
\midrule
S1 & LLM & ``take the red cube to dispatch''---single-object delivery, closed-loop recovery \\
S2 & Baseline & same task, hardcoded---deterministic reference \\
S3 & LLM & ``move the red cube to dispatch, then the green cube''---ordered multi-step \\
S5 & LLM & ``tidy up the warehouse''---ambiguous command interpretation \\
\bottomrule
\end{tabular}
\end{table}

\paragraph{Reporting policy (development vs.\ frozen system)}
The system was built and debugged incrementally, and the raw logs include many
early failures produced \emph{while core reliability fixes were still being
written} (e.g.\ before the approach-nudge margin and a sign bug in the reverse
maneuver were corrected; Section~\ref{sec:lessons}). We do not report those as
system performance. All results below are for the \emph{frozen configuration}
(runs from 2026-06-25 onward), and we report the true success rate of that window,
including runs that still failed on the frozen system.

% ============================================================================
\section{Results}\label{sec:results}
% ============================================================================
\begin{table}[t]
\caption{Frozen-system results. LLM successes typically include $\geq1$ autonomous
replan. Baseline measured on its tuned configuration.}
\label{tab:results}
\centering
\small
\begin{tabular}{@{}llccl@{}}
\toprule
Scenario & System & $N$ & Success & Mean time (success) \\
\midrule
S1 single   & LLM      & 16 & \textbf{81\%} (13/16) & 454\,s (232--714) \\
S2 single   & Baseline & 5  & 100\% (5/5)\,$^\dagger$ & 194\,s (143--249) \\
S3 two-cube & LLM      & 4  & 50\% (2/4)   & 596\,s (545--648) \\
S5 tidy-up  & LLM      & 1  & qualitative  & --- \\
\bottomrule
\end{tabular}
\\[2pt]
\raggedright\footnotesize $^\dagger$Earlier baseline failures in the raw log
predate the reliability fixes and are excluded as development runs; on the tuned
configuration the baseline is deterministic.
\end{table}

\subsection{S1: single-object delivery and autonomous recovery}
On the frozen system the LLM succeeded in $13/16$ ($81\%$) single-object
deliveries, taking a mean of $454$~s ($232$--$714$~s) and $5$--$8$ LLM calls per
success. The key qualitative finding is that \emph{nearly every success includes
an autonomous recovery}: the LLM's first \texttt{pick} from the navigation
standoff fails ``out of reach,'' and---reading the error---it issues
\texttt{drive\_distance(0.30)} and re-\texttt{pick}s successfully, without any
task-specific scripting. This is precisely the behavior the fixed baseline lacks.

The three residual failures were \emph{not} mechanical: the arm demonstrably
reaches the cube on fifteen other runs. They were \emph{planner loops}, in which
the LLM exhausted the turn/replan budget (e.g.\ $26$ calls / $7$ replans) while
oscillating between corrective moves. We report this as the characteristic failure
mode of a flexible planner.

\subsection{S2: deterministic baseline}
On its tuned configuration the baseline completed $5/5$ single-object deliveries
in $143$--$249$~s. Being perception-blind and branch-free, it is fast and
repeatable, but it cannot react: it succeeds precisely when the world matches the
program it was given, and by construction has no recovery path when it does not.
This is the intended contrast---determinism versus adaptability---rather than a
raw win/loss.

\subsection{S3: ordered multi-object sequences}
On two-cube ordered delivery (``red first, then green''), the LLM succeeded in
$2/4$ runs. Successes were clean---$16$ LLM calls, correct ordering, both cubes
delivered---but both failures hit the turn limit ($31$ calls, $7$--$10$ replans).
The longer the task, the more opportunities the planner has to enter a replan
loop; reliability therefore \emph{degrades with sequence length}. We treat this as
a measured scalability limitation, not a headline result.

\subsection{S5: ambiguous command interpretation}
Given ``tidy up the warehouse,'' the LLM \emph{correctly interpreted} the
ambiguity---it inferred that tidying meant relocating cubes from the table to the
dispatch zone, navigated to the table, and began attempting picks. Execution then
failed: after the approach nudge, a detection cycle returned no objects, starving
the subsequent \texttt{pick}; the planner then flailed and incorrectly concluded
that the arm's reach was insufficient (contradicted by S1/S3). The root cause
appears to be a perception field-of-view gap after the local base move, not a
reachability wall. S5 thus demonstrates that the LLM \emph{understands} ambiguous
instructions while surfacing a concrete perception-robustness limitation; the
baseline cannot attempt S5 at all, as it performs no language interpretation.

% ============================================================================
\section{Discussion: Simulation--Robotics Integration Lessons}\label{sec:lessons}
% ============================================================================
Much of the engineering effort---and much of the transferable value---lay in
physical-interaction bugs that a pure planning study would never surface:

\begin{itemize}
\item \textbf{Unsigned reverse distance.} The post-pick reverse measured distance
as $\sqrt{dx^2+dy^2}$, which is always positive. With the arm extended, its
inertia could pull the base \emph{forward}, yet this was counted as progress
backward, parking the robot against the table. The fix is a signed projection onto
the backward heading.
\item \textbf{Reach margin.} A $0.15$~m approach nudge was insufficient for
pre-grasp IK; $0.30$~m was required. A small workspace-geometry constant dominated
overall reliability.
\item \textbf{Grasp impulse corrupts localization.} The gripper attach impulse
jolted the chassis, injecting phantom LiDAR returns and corrupting the pose graph.
We therefore ran experiments on a frozen static map---which in turn means our
results do not include real localization drift, a limitation revisited below.
\item \textbf{Ordering and physics.} Reversing with the arm extended dragged the
base; moving the arm home first (at reduced velocity) fixed it. A single gripper
joint oscillated at its limit until velocity scaling was halved.
\item \textbf{Framework foot-guns.} Subclassing the rclpy \texttt{Node} and
reusing its internal \texttt{\_clients} attribute crashed the planner; navigation
to the far zone exceeded a default service timeout and had to be raised to
$300$~s.
\end{itemize}

% ============================================================================
\section{Limitations and Threats to Validity}
% ============================================================================
This is a proof-of-concept study in simulation only, with no real-robot transfer.
Sample sizes are small ($N=16/5/4/1$). We evaluate a single LLM (Claude Haiku~4.5)
and do not compare across models. A frozen static map removes the localization
drift a deployed system would face---and which our longer runs began to
expose---motivating tighter sensor fusion and periodic relocalization as future
work. Finally, baseline and LLM success rates are measured over different
denominators and reflect different configurations; we therefore compare them
qualitatively (determinism vs.\ adaptability) rather than as a single head-to-head
percentage.

% ============================================================================
\section{Conclusion}
% ============================================================================
A general-purpose LLM, constrained to a fixed tool interface and closed with a
structured scene state, can drive a complete ROS~2 mobile-manipulation stack from
a natural-language sentence and---critically---recover autonomously from realistic
execution failures, succeeding on $81\%$ of single-object deliveries where a
hardcoded baseline is fast but rigid. Reliability degrades on longer sequences and
under ambiguous commands whose execution stresses perception, and we report these
limits candidly alongside the successes. We believe the honest failure catalogue
and the concrete integration lessons are as valuable to practitioners as the
headline success rate. Future work includes real-robot transfer, sensor fusion to
support long-horizon runs, cross-model comparison, and a controlled
disturbance-recovery experiment that the perception-driven planner is
architecturally positioned to win.

% ============================================================================
\section*{Reproducibility}
% ============================================================================
All packages, the tool schemas, the LLM system prompt, the world file, and the
experiment/logging code are in the project repository. The companion technical
dossier documents every version, coordinate, tool, and result with provenance.

\begin{thebibliography}{1}
\bibitem{saycan} M.~Ahn \emph{et al.}, ``Do As I Can, Not As I Say: Grounding
Language in Robotic Affordances,'' 2022.
\bibitem{ros2} S.~Macenski \emph{et al.}, ``Robot Operating System 2: Design,
architecture, and uses in the wild,'' \emph{Science Robotics}, 2022.
\bibitem{nav2} S.~Macenski \emph{et al.}, ``The Marathon 2: A Navigation System
(Nav2),'' IROS, 2020.
\bibitem{moveit} D.~Coleman \emph{et al.}, ``Reducing the Barrier to Entry of
Complex Robotic Software: a MoveIt! Case Study,'' 2014.
\bibitem{yolo} G.~Jocher \emph{et al.}, ``Ultralytics YOLOv8,'' 2023.
% Add: LLM tool-use / code-as-policies / inner-monologue references for camera-ready.
\end{thebibliography}

\end{document}
